\scanpage{74}
% lines 1–5, linear extension
On peut aussi considérer l'extension linéaire \unsure{de|des} $\pi_i$ :
\[
E_{\unsure{j|i}}\longrightarrow L_i
\]
et la condition précédente signifie simplement que
\[
\pi_i(\deleted{\unsure{x}}u_i)=\unsure{\chi|\gamma}\cdot u_i\qquad\text{si }u_i\in L_i\subset E_i.
\]
% lines 6–10, fundamental formulas and right-hand glosses
Ceci posé, on aura les formules fondamentales qui suivent.

\textcircled{1} \uuline{bilinéarité} de $x_1*x_2$,
\margin{i.e. correspond à \uline{appl. lin.} $\boxed{E_1\otimes E_2\to E}$}
\[
\boxed{\psi(x_1*x_2)=\psi_1(x_1)\psi_{\unsure{1|2}}(x_2)}\tag*{\textcircled{2}}
\]
\margin{compatibilité aux augmentations}
où
\[
\begin{aligned}
E_1&\overset{\psi_1}{\mathrel{\text{---}}}\underline O\\
E_2&\overset{\psi_2}{\mathrel{\text{---}}}\underline O\\
E=E_1*E_2&\xrightarrow{\psi}\underline O
\end{aligned}
\]
% lines 11–15, augmentation fibres
sont les augmentations. En particulier, $*$ induit
\[
\begin{aligned}
P_1\times P_2&\longrightarrow P&\quad&\text{où}&P_i&=\psi_i^{-1}(\{1\})\\
x_1,x_2&\longmapsto x_1*x_2&&&P&=\psi^{-1}(\{1\})\\[-2pt]
&\text{induit par }*
\end{aligned}
\]
% lines 16–20, translation formulas, left annotation
\[
\boxed{\begin{aligned}
(x_1+u_1)*x_2&=x_1*x_2+\deleted{\unsure{\pi_2}}u_1\otimes\pi_2(x_2)\\
x_1*(x_2+u_2)&=x_1*x_2+\pi_1(x_1)\otimes u_2
\end{aligned}}
\]
\[
x_i\in\Gamma E_{\unsure{j|i}},\qquad u_i\in\Gamma L_i
\]
% slanted
\margin{\textcircled{3} Ces formules \unsure{utiles} sont celles \unsure{pour} \unsure{surtout} en se \unsure{restreignant} à des arguments $x_1\in\Gamma P_1,\ x_2\in\Gamma P_2$}
d'où aussitôt
\[
\boxed{\begin{aligned}
(x_1+u_1)*(x_2+u_2)={}&x_1*x_2+u_1\otimes\pi_2(x_2)+\pi_1(x_{\unsure{1|0}})\otimes u_2\\
&+\chi\,u_1\otimes u_2
\end{aligned}}
\]
% lines 21–26, restrictions at foot
Ces formules résultent de la bilinéarité $+$
\[
\boxed{\begin{aligned}
u_1*x_2&=u_1\otimes\pi_2(x_2)\\
x_1*u_2&=\pi_1(x_1)\otimes u_2
\end{aligned}}
\]
\deleted{ce qui} dont chacune implique (en faisant $x_2=u_2$ resp $x_1=u_1$)
\[
\boxed{u_1*u_2=\chi\,u_1\otimes u_2},
\]
compte tenu de $\pi_1(u_1)=\chi u_1,\quad\pi_2(u_2)=\chi(u_2)$
% archive sheet number: 73
